\hypertarget{cont1:paper-ix-finite-polynomial-lifting-and-shared-corrections}{%
\subsection{Chapter IX: finite polynomial lifting and shared
corrections}\label{cont1:paper-ix-finite-polynomial-lifting-and-shared-corrections}}

The nonlinear Taylor argument supplies the correction formula used in
one-step polynomial lifting. The statements concern a fixed
finite list of integer polynomials. They hold for every positive integer
modulus, including composite moduli and modulus one. No
Jacobian-surjectivity assumption is needed.

\hypertarget{cont1:statement-and-notation}{%
\subsubsection{Statement and notation}\label{cont1:statement-and-notation}}

Let \[
F=(F_1,\ldots,F_e)\in\mathbb Z[X_1,\ldots,X_d]^e,
\quad q\ge1,\quad L\ge1,
\quad x\in\mathbb Z^d,
\quad F(x)\equiv0\pmod{q^L}.
\] Write \(R_q=\mathbb Z/q\mathbb Z\), let \(\rho_q\) denote
coordinatewise reduction, and put \[
r_x=F(x)/q^L\in\mathbb Z^e,
\qquad A_x=DF(x):\mathbb Z^d\longrightarrow\mathbb Z^e,
\qquad J_x=\rho_q\circ A_x.
\] Here \(A_x\) is computed from the actual formal partial derivatives
of the input polynomials. The Lean map \(J_x\) has domain
\(\mathbb Z^d\) and codomain \(R_q^e\), viewed as integer modules. It
factors through reduction to \(R_q^d\); the cokernel below is presented
directly using its integer-domain image.

Define the finite obstruction module and its distinguished class by \[
\mathcal O_x=R_q^e/\operatorname{im}J_x,
\qquad o_x=[\rho_q(r_x)]\in\mathcal O_x.
\] The checked criterion is \[
\boxed{
 o_x=0
 \iff
 \exists y\in\mathbb Z^d:\quad
 y\equiv x\pmod{q^L},\qquad F(y)\equiv0\pmod{q^{L+1}}.
}
\] The residual exists by the lower congruence and is unique because
\(q>0\). The quotient is finite since its ambient module \(R_q^e\) is
finite. The theorem permits zero variables or zero equations as well.

\hypertarget{cont1:the-checked-taylor-bridge}{%
\subsubsection{The checked Taylor
bridge}\label{cont1:the-checked-taylor-bridge}}

For every \(P\in\mathbb Z[X_1,\ldots,X_d]\), integer \(a\), and
\(x,t\in\mathbb Z^d\), Lean proves \[
\boxed{
 a^2\mid P(x+at)-P(x)-a\sum_{i=1}^{d}(\partial_iP)(x)t_i.
}
\] The proof uses polynomial induction: constants, addition, and
multiplication by a variable. The product rule for the actual partial
derivatives controls the multiplication step. It does not assume a
Taylor identity or a liftability condition.

Taking \(a=q^L\), the inequality \(2L\ge L+1\) gives \[
F(x+q^Lt)\equiv q^L(r_x+A_xt)\pmod{q^{L+1}}.
\] Since \(q^L\ne0\) in the integers, cancellation in divisibility
yields \[
F(x+q^Lt)\equiv0\pmod{q^{L+1}}
\iff \rho_q(r_x)+J_xt=0.
\] Every integer representative of the lower residue has the form
\(x+q^Lt\), and this correction equation has a solution exactly when
\(o_x=0\).

The positive-depth hypothesis is necessary for the general nonlinear
theorem. For \(F(X)=X^2-1\), \(q=2\), \(L=0\), \(x=0\), the Jacobian is
zero and the residual class is nonzero, yet \(y=1\) is a solution modulo
two in the same residue modulo one. Affine systems retain the previously
checked \(L\ge0\) theorem.

\hypertarget{cont1:independence-of-the-integer-representative}{%
\subsubsection{Independence of the integer
representative}\label{cont1:independence-of-the-integer-representative}}

Suppose \(x'=x+q^Ls\) represents the same lower residue, with exact
residual \(r_{x'}\). Polynomial evaluation respects coordinatewise
congruences, including when applied to each partial derivative. Hence \[
J_{x'}=J_x,
\qquad
\rho_q(r_{x'})=\rho_q(r_x)+J_xs.
\] The two cokernels therefore have the same ambient module and the same
defining submodule, and \[
[\rho_q(r_{x'})]=[\rho_q(r_x)].
\] Lean states this using one explicit common reduced map, avoiding an
implicit comparison between differently indexed quotient types. This
establishes independence of the integer representative within the fixed
polynomial presentation.

Changing the polynomial presentation is a separate question. Appending a
zero equation can change the raw cokernel by adding a copy of \(R_q\),
even though it does not change the solutions.

\hypertarget{cont1:two-polynomial-blocks-and-one-shared-lift}{%
\subsubsection{Two polynomial blocks and one shared
lift}\label{cont1:two-polynomial-blocks-and-one-shared-lift}}

Let \(G\in\mathbb Z[X_1,\ldots,X_d]^f\) be a second fixed system with
the same variables and the same lower root \(x\). Let
\(K_x=\rho_q\circ DG(x)\), and define \[
\mathcal C_x=\mathbb Z^d/(\ker J_x+\ker K_x).
\] This quotient is finite: it is a quotient of the finite image of
\(J_x\), or equivalently of the finite image of \(K_x\).

Suppose \(u,v\in\mathbb Z^d\) are separately valid corrections: \[
F(x+q^Lu)\equiv0\pmod{q^{L+1}},
\qquad
G(x+q^Lv)\equiv0\pmod{q^{L+1}}.
\] Then \[
\boxed{
\begin{gathered}
[u]=[v]\text{ in }\mathcal C_x
\iff
\exists y\equiv x\pmod{q^L}:\\
F(y)\equiv0\pmod{q^{L+1}},\quad G(y)\equiv0\pmod{q^{L+1}}.
\end{gathered}
}
\] The labels are independent of the choices of valid \(u\) and \(v\).
An unconditional version asks for separately valid corrections with
equal labels; it does not assume that each block is liftable.

Taylor transport turns a common correction \(t\) into the simultaneous
equations \(J_xt=J_xu\) and \(K_xt=K_xv\). The previously checked
common-quotient theorem gives precisely the displayed equivalence.
Separate solvability alone does not imply joint solvability; the
controls include \(F(X)=X\), \(G(X)=X-2\), starting at \(x=0\) modulo
two and lifting to modulus four.
